\subsection{Legacy code comprehension and specification recovery}
\label{subsec:rw8}

Most software in production is not written from scratch but inherited, and the reference pipeline proposed in this article must therefore confront brownfield codebases in which intent is implicit, distributed, and partly erroneous. The disciplines of legacy code comprehension and specification recovery have a long history of grappling with precisely this condition, and they supply both the conceptual vocabulary and the technical building blocks on which the present work draws. We group the relevant literature into four strands: the contractual stance that motivates recovering machine-checkable obligations at all; behaviour-preservation and characterisation practices for safely manipulating opaque code; dynamic and static specification mining; and the recent turn to learned, language-model--based recovery of formal contracts.

The contractual stance originates with Meyer's Design by Contract \cite{meyer1992dbc}, which recast a software component's interface as a set of mutual obligations---preconditions, postconditions, and invariants---between caller and callee. This framing is foundational to the present article: the contract layer it proposes is exactly a set of such obligations, and the compositional propagation it inherits from Article~3 of this programme \cite{edmonds_article3} is a direct consequence of contracts composing across method boundaries. Meyer assumed contracts would be authored by the developer up front, however, whereas the brownfield reality is that the contract must be \emph{recovered} from code that predates any contractual discipline. The terminology for that recovery activity was rationalized by Chikofsky and Cross \cite{chikofsky1990reverse}, whose taxonomy distinguishes reverse engineering, design recovery, and reengineering; their notion of \emph{design recovery}---reconstructing higher-level abstractions from artefacts and domain knowledge beyond the code itself---anticipates the multi-source triangulation (code, tickets, documentation) that this article's startup phase performs.

The behaviour-preservation strand is exemplified by Feathers' influential treatment of legacy systems \cite{feathers2004legacy}, which defines legacy code operationally as code lacking tests and introduces \emph{characterisation tests}: tests that pin down what the code currently does rather than what it should do, establishing a regression baseline before any change. This is conceptually identical to what we call the characterisation phase, and to the precise claim we make about what verification against a code-derived specification actually proves---namely a faithful description of current behaviour, not its correctness. The crucial difference is one of medium and coverage: a characterisation test fixes behaviour at sampled inputs, whereas a recovered formal contract checked by Extended Static Checking (ESC) constrains behaviour over \emph{all} inputs, and disagreements between the recovered contract and independently sourced intent become candidate bugs rather than silently captured behaviour.

Specification mining sought to automate this recovery. Ammons et~al.\ \cite{ammons2002mining} coined the term and inferred finite-state API-protocol specifications from execution traces under the assumption that common behaviour is usually correct behaviour. In parallel, Ernst and colleagues' Daikon system \cite{ernst2007daikon} dynamically detected likely value invariants---ranges, ordering, linear relationships---over observed runs, and Nimmer and Ernst \cite{nimmer2001static} integrated Daikon with ESC/Java so that dynamically proposed invariants could be statically confirmed. This propose-then-verify pairing is the closest antecedent to our pipeline's structure, yet it differs in two decisive ways. First, dynamic mining is bounded by test coverage and produces invariants that may be accidental, whereas our drafts are inferred statically from the abstract syntax tree by the JML-Inferrer and discharged by OpenJML over all inputs. Second, and more fundamentally, neither Daikon nor classical specification mining addresses \emph{provenance}: a mined invariant is an undifferentiated guess, with no record of whether it reflects ratified human intent, documented behaviour, or merely incidental regularity in the traces.

The recent language-model strand attacks the natural-language gap that mining never could. Endres et~al.\ \cite{endres2024nl2postcond} translate informal intent into formal postconditions capable of catching real Defects4J bugs; SpecGen \cite{ma2025specgen} generates verifiable specifications through conversational prompting and mutation, outperforming Daikon and Houdini; and a broader line of work probes contract synthesis at scale \cite{lahiri2025nl2contract}. These results validate the feasibility of the elaboration layer above this article's pivot, but they treat specification generation as a one-shot, self-graded translation, leaving open the very tensions this article foregrounds: the contract is taken as trusted output rather than ratified against pinned behavioural examples, the generating model is not held independent of the implementer, and a single assurance level is assumed rather than a tiered degradation that records what could only be tested or bounded-checked.

The gap this article addresses is thus the synthesis these strands have not achieved. Prior recovery work mines or generates specifications, but it does not (i) reconcile code-, ticket-, and documentation-derived intent into a single provenance- and assurance-tagged contract whose internal \emph{disagreements} are surfaced as candidate bugs; (ii) treat the recovered contract as a non-trusted artefact that drives a Socratic, example-based ratification dialogue rather than gating on zero errors; (iii) enforce an independent-authority separation between the agent that elaborates the contract and the agent that satisfies it; or (iv) promote compositional contract propagation \cite{edmonds_article3} into a lifecycle event that flags which callers a library or dependency change breaks behaviourally without executing them. The reference pipeline integrates legacy comprehension, machine-checkable contracts, and change-impact analysis into one verification-centric loop, with the recovered contract---not the disposable code---as the durable artefact.
